\section{Metal Rendering Architecture}
\label{sec:metal}

\subsection{Device Objects and Lifetime}

A single \texttt{MTLDevice}, one \texttt{MTLCommandQueue} for rendering, and a
second low-priority queue for background work (thumbnails, export) are created
at application launch and live for the process lifetime. The render queue is
owned exclusively by the render actor of Section~\ref{sec:arch}; no other
context may enqueue on it.

All shader functions live in a single default \texttt{MTLLibrary} compiled at
build time. Pipeline states are compiled lazily, cached by a key comprising the
function name and the function-constant values, and persisted to disk through
the Metal binary archive so that second launch avoids recompilation:

\begin{lstlisting}[language=Swift, caption={Pipeline state cache with binary archive.}]
actor PipelineCache {
    private var states: [PipelineKey: MTLComputePipelineState] = [:]
    private let archive: MTLBinaryArchive?

    func state(for key: PipelineKey) throws -> MTLComputePipelineState {
        if let s = states[key] { return s }
        let desc = MTLComputePipelineDescriptor()
        desc.computeFunction = try library.makeFunction(
            name: key.function, constantValues: key.constants)
        desc.label = key.debugLabel
        // Prefer the archived binary; fall back to JIT and record it.
        if let a = archive { desc.binaryArchives = [a] }
        let s = try device.makeComputePipelineState(
            descriptor: desc, options: [], reflection: nil).0
        states[key] = s
        return s
    }
}
\end{lstlisting}

Cold-start pipeline compilation for the full graph is approximately 180\,ms on
an A15, which would violate NFR-04. The binary archive reduces this to under
20\,ms on subsequent launches, and the first launch compiles only the subset of
pipelines the capture path needs, deferring the develop-path pipelines to a
background task started after the camera is live.

\subsection{Precision Policy}
\label{sec:precision}

Constraint CR-05 forbids \texttt{half} for scene-referred linear data. The full
policy is:

\begin{table}[htbp]
\caption{Precision Policy by Pipeline Domain}
\label{tab:precision}
\begin{center}
\renewcommand{\arraystretch}{1.15}
\begin{tabularx}{\columnwidth}{@{}ll>{\raggedright\arraybackslash}X@{}}
\toprule
\textbf{Domain} & \textbf{Format} & \textbf{Rationale} \\
\midrule
Scene linear & \texttt{rgba32f} & 18-stop range; \texttt{half} subnormals quantize shadows. \\
Log exposure & \texttt{rgba32f} & Feeds \texttt{exp}; error amplifies. \\
Halation source & \texttt{r16f} & Single channel, large values, low precision need. \\
Halation pyramid & \texttt{rgba16f} & Blurred data; error is spatially smoothed. \\
Negative density & \texttt{rgba16f} & Range $[0,4]$, needs $10^{-3}$; \texttt{half} gives $5\times10^{-4}$. \\
Grain field & \texttt{rgba16f} & Zero-mean unit-variance. \\
Print density & \texttt{rgba16f} & As negative density. \\
Display linear & \texttt{rgba16f} & Pre-encode. \\
Preview output & \texttt{bgra8unorm\_srgb} & Direct to drawable. \\
Export output & \texttt{rgba16unorm} & 16-bit TIFF/PNG path. \\
\bottomrule
\end{tabularx}
\end{center}
\end{table}

Arithmetic within kernels is performed in \texttt{float} regardless of storage
format \cite{metalspec}. The \texttt{half} storage saves bandwidth, which is the
binding constraint on a mobile GPU \cite{metalbest}; it does not save ALU, and
using \texttt{half} arithmetic in the density model measurably degrades the toe.

\subsection{Resource Management}

Textures are allocated from a small number of \texttt{MTLHeap} objects with
manual aliasing. The render graph (Section~\ref{sec:rendergraph}) computes the
lifetime interval of every intermediate texture, then performs a linear-scan
allocation so that textures with disjoint lifetimes share memory.

For the full pipeline at 48\,MP, the naive allocation is 14 intermediates of
$4032 \times 3024 \times 8$ bytes $= 97.5$\,MB each, totalling 1.37\,GB and
violating NFR-03. Lifetime aliasing reduces the live set to four full-resolution
intermediates plus the pyramid, totalling 442\,MB. The remaining headroom
absorbs the source RAW and the export buffer.

\begin{lstlisting}[language=Swift, caption={Heap-backed transient texture allocation.}]
final class TransientPool {
    private let heap: MTLHeap
    private var live: [ResourceID: MTLTexture] = [:]

    func acquire(_ desc: MTLTextureDescriptor, id: ResourceID) -> MTLTexture {
        desc.storageMode = .private
        desc.hazardTrackingMode = .untracked   // graph inserts barriers
        let tex = heap.makeTexture(descriptor: desc)!
        live[id] = tex
        return tex
    }
    /// Called when the graph determines the resource's last reader has run.
    func release(_ id: ResourceID) {
        live[id]?.makeAliasable()
        live[id] = nil
    }
}
\end{lstlisting}

Hazard tracking is disabled on heap resources and barriers are inserted
explicitly by the graph from the dependency edges it already knows. Automatic
hazard tracking on an aliased heap is conservative to the point of serializing
the entire graph.

\subsection{Threadgroup Sizing}

All compute kernels are dispatched with \texttt{dispatchThreads} (non-uniform
threadgroups), which is supported on all target devices and removes the bounds
check from the common case. Threadgroup size is chosen per pipeline from the
device's reported \texttt{maxTotalThreadsPerThreadgroup} and
\texttt{threadExecutionWidth}:

\begin{equation}
w = \mathtt{threadExecutionWidth},\quad h = \left\lfloor \frac{\mathtt{maxTotal}}{w} \right\rfloor,
\end{equation}

clamped to $h \leq 8$ for pointwise kernels (which are bandwidth-bound and
benefit from wide, shallow groups with good cache locality) and $h = w$ for
separable filters (which benefit from square tiles when using threadgroup
memory). The values are queried, not hardcoded: on Apple silicon $w = 32$, but
assuming it is a correctness hazard on any future device.

\subsection{Kernel Fusion and Tile Memory}

Each pass that reads and writes a full-resolution texture costs, at 48\,MP and
\texttt{rgba16f}, 97.5\,MB read plus 97.5\,MB write. At a realistic 60\,GB/s of
achievable bandwidth that is 3.25\,ms per pass in memory traffic alone. With
fourteen passes the pipeline would be bandwidth-bound at 45\,ms before any
arithmetic --- a factor that dominates every other consideration.

Fusion is therefore the primary optimization, and the graph performs it
automatically for chains of pointwise passes (Section~\ref{sec:rendergraph}).
Beyond that, two Metal-specific mechanisms are used:

\paragraph{Imageblock (tile) memory} On Apple GPUs, a render pass can retain
per-tile data in on-chip imageblock storage across multiple draw or dispatch
calls without a round trip to device memory. The optical warp, the pointwise
develop chain, and the output transform are expressed as tile dispatches within
a single render pass, so the intermediate never leaves tile memory. This
eliminates six full-resolution round trips.

\paragraph{Memoryless attachments} Any attachment that is written and consumed
entirely within one render pass is declared \texttt{.memoryless}, so it is never
backed by device memory at all.

\begin{lstlisting}[language=Metal, caption={Fused pointwise develop chain as a tile dispatch.}]
struct DevelopTile {
    half4 color;
};

[[kernel]] void develop_tile(
    imageblock<DevelopTile, imageblock_layout_implicit> blk,
    constant DevelopParams& P    [[buffer(0)]],
    texture3d<half, access::sample> lut [[texture(0)]],
    ushort2 lid [[thread_position_in_threadgroup]])
{
    threadgroup_imageblock DevelopTile* px = blk.data(lid);
    float3 e = float3(px->color.rgb);

    e = apply_vignette_and_warp_correction(e, P);   // fused optics
    float3 x = log10(max(e, P.epsFloor));           // to log exposure
    float3 D = sample_baked_lut(lut, x, P);         // pointwise chain
    px->color = half4(half3(D), px->color.a);
}
\end{lstlisting}

\subsection{Pointwise Chain Baking}
\label{sec:lutbake}

The pointwise portion of the pipeline --- layer balance, characteristic curve,
mask depletion, printing density matrix, print curve, silver retention, neutral
axis, and output transform --- is a fixed function $\Rspace^3 \to \Rspace^3$ for
a given parameter set. Evaluating it per pixel costs roughly 140 ALU operations
including six transcendentals. Baking it into a 3D lookup table reduces that to
one texture fetch with hardware trilinear filtering, or eight fetches with
software tetrahedral interpolation.

\paragraph{Shaper} The input domain is log exposure over $[-4, +2]$, which is
uniform in the perceptually relevant variable, so the shaper is simply the
affine map

\begin{equation}
s(x) = \frac{x + 4}{6} \in [0,1].
\end{equation}

Uniform sampling in $\logE$ is essential; a LUT uniform in linear exposure would
allocate almost all of its resolution to highlights and band the shadows badly.

\paragraph{Grid size} A $45^3$ grid is used. The error analysis is
straightforward: the maximum second derivative of the composite transfer occurs
at the print stock's shoulder, where $|D''| \lesssim 9$ per unit of $\logE$.
Trilinear interpolation error is bounded by $\tfrac{1}{8}|D''|\Delta^2$ with
$\Delta = 6/44 = 0.136$, giving $2.1 \times 10^{-2}$ density --- too large.
Tetrahedral interpolation on the same grid reduces the error by approximately a
factor of three for this class of function, and the residual $7\times10^{-3}$
density is below the $10^{-2}$ visibility threshold. A $65^3$ grid would halve
the error again at $2.7\times$ the memory; $45^3 \times 8$ bytes is 729\,KB,
which fits comfortably and is rebuilt in under 1\,ms.

\paragraph{Rebuild policy} The LUT is rebuilt on the GPU by a
$45\times45\times45$ dispatch evaluating the analytic chain --- 91{,}125 threads,
a trivial dispatch. It is invalidated by any change to a pointwise parameter and
is not invalidated by changes to grain, halation, interlayer, or optical
parameters. The dependency tracking that determines this is specified in
Section~\ref{sec:rendergraph}.

\paragraph{Tetrahedral interpolation} Trilinear interpolation of a 3D LUT is
not exact on the neutral axis: the eight corners of a cell do not in general
interpolate to a neutral result, producing faint colored artifacts in gray
gradients. Tetrahedral interpolation \cite{kang2006} decomposes the cube into
six tetrahedra
along the main diagonal, which \emph{is} the neutral axis, so neutrals are
reproduced exactly.

\begin{lstlisting}[language=Metal, caption={Tetrahedral 3D LUT interpolation.}, label={lst:tetra}]
inline float3 tetrahedral(texture3d<float, access::read> lut,
                          float3 uvw, uint N)
{
    float3 p = uvw * float(N - 1);
    uint3  i = uint3(floor(p));
    i = min(i, uint3(N - 2));
    float3 f = p - float3(i);

    // Eight cell corners
    float3 c000 = lut.read(i + uint3(0,0,0)).rgb;
    float3 c111 = lut.read(i + uint3(1,1,1)).rgb;

    // Select the tetrahedron by ordering of (fx, fy, fz)
    float3 v1, v2;  float w1, w2, w3;
    if (f.x > f.y) {
        if (f.y > f.z)      { v1 = lut.read(i+uint3(1,0,0)).rgb;
                              v2 = lut.read(i+uint3(1,1,0)).rgb;
                              w1 = f.x-f.y; w2 = f.y-f.z; w3 = f.z; }
        else if (f.x > f.z) { v1 = lut.read(i+uint3(1,0,0)).rgb;
                              v2 = lut.read(i+uint3(1,0,1)).rgb;
                              w1 = f.x-f.z; w2 = f.z-f.y; w3 = f.y; }
        else                { v1 = lut.read(i+uint3(0,0,1)).rgb;
                              v2 = lut.read(i+uint3(1,0,1)).rgb;
                              w1 = f.z-f.x; w2 = f.x-f.y; w3 = f.y; }
    } else {
        if (f.z > f.y)      { v1 = lut.read(i+uint3(0,0,1)).rgb;
                              v2 = lut.read(i+uint3(0,1,1)).rgb;
                              w1 = f.z-f.y; w2 = f.y-f.x; w3 = f.x; }
        else if (f.z > f.x) { v1 = lut.read(i+uint3(0,1,0)).rgb;
                              v2 = lut.read(i+uint3(0,1,1)).rgb;
                              w1 = f.y-f.z; w2 = f.z-f.x; w3 = f.x; }
        else                { v1 = lut.read(i+uint3(0,1,0)).rgb;
                              v2 = lut.read(i+uint3(1,1,0)).rgb;
                              w1 = f.y-f.x; w2 = f.x-f.z; w3 = f.z; }
    }
    float w0 = 1.0f - w1 - w2 - w3;
    return w0*c000 + w1*v1 + w2*v2 + w3*c111;
}
\end{lstlisting}

\subsection{Separable Filters}

Gaussian blurs at the small radii used inside the pyramid (Section
\ref{sec:recursive-blur}) are implemented as separable two-pass convolutions with
threadgroup-memory staging. For a $\sigma \le 3$ kernel the tap count is 13, and
linear-sampling pairing halves the fetches to 7.

\texttt{MPSImageGaussianBlur} is deliberately \emph{not} used. It is
well-optimized, but it clamps at the texture edge in a way that is not
configurable, it allocates internally, and it cannot be fused into a tile
dispatch. The application-owned implementation is within 8\% of its throughput
and composes with the rest of the graph.

\subsection{Function Constants and Specialization}
\label{sec:specialization}

Branching on parameter state inside a kernel costs both instructions and
register pressure. Instead, boolean and small-integer configuration is exposed
as Metal function constants, and the pipeline cache keys on their values, so
each configuration compiles to a specialized kernel with the branches removed:

\begin{lstlisting}[language=Metal, caption={Function-constant specialization.}]
constant bool  kMonochrome        [[function_constant(0)]];
constant bool  kSilverRetention   [[function_constant(1)]];
constant bool  kSpatialActivity   [[function_constant(2)]];  // chemical streaking
constant uint  kPrintStockModel   [[function_constant(3)]];

kernel void develop_pointwise(/* ... */) {
    float3 D = characteristic_curve(x, negParams);
    if (kSpatialActivity) {                 // compiled out when false
        D = apply_spatial_activity(D, x, streakTex, gid);
    }
    // ...
}
\end{lstlisting}

The number of live specializations is bounded: three booleans and one four-value
enum give 32 combinations, of which fewer than 10 occur in practice. The cache
holds all of them.

\subsection{Command Buffer Structure}

A preview render is one command buffer containing, in order: an optional LUT
rebuild dispatch; a render pass containing the fused tile chain; two compute
encoders for the spatial stages; and the final blit or drawable present. The
buffer is scheduled with a completion handler that signals the render actor,
which then picks up any pending recipe.

Export uses a separate command buffer on the same queue, submitted with
\texttt{MTLCommandBufferErrorOption.\allowbreak{}encoderExecutionStatus} so that
a GPU fault
during a long export is attributable to a specific encoder rather than crashing
opaquely.

\subsection{Debugging and Instrumentation}

Every encoder is labelled, every resource is labelled, and debug groups wrap
each logical stage, so a Metal frame capture reads as the pipeline diagram of
Figure~\ref{fig:pipeline}. In debug builds, Metal shader validation and
\texttt{MTLCaptureManager} programmatic capture on a triggered gesture are
enabled. A build-time flag inserts a per-stage GPU timestamp using
\texttt{MTLCounterSampleBuffer}, and the resulting per-stage timings are
surfaced in a developer overlay --- which is how the numbers in
Section~\ref{sec:performance} are collected.
